\documentclass[12pt,oneside]{book}

% =========================
% PAKET
% =========================
\usepackage[a4paper,
top=3cm,
bottom=3cm,
left=3.5cm,
right=3.5cm]{geometry}

\usepackage{setspace}
\usepackage{lmodern}
\usepackage[T1]{fontenc}
\usepackage{microtype}
\usepackage{titlesec}
\usepackage{fancyhdr}
\usepackage{emptypage}
\usepackage{tocloft}
\usepackage{hyperref}

% =========================
% BAHASA
% =========================
\renewcommand{\contentsname}{Daftar Isi}

% =========================
% SPASI BARIS
% =========================
\onehalfspacing

% =========================
% HEADER & FOOTER
% =========================
\pagestyle{fancy}
\fancyhf{}
\fancyfoot[C]{\thepage}
\renewcommand{\headrulewidth}{0pt}

% =========================
% GAYA JUDUL BAB
% =========================
\titleformat{\chapter}
{\normalfont\Huge\bfseries\centering}
{}
{0pt}
{\Huge}

\titlespacing*{\chapter}{0pt}{-10pt}{40pt}

% =========================
% DOKUMEN
% =========================
\begin{document}
	
	% =========================
	% TITLE PAGE (TANPA NOMOR)
	% =========================
	\pagenumbering{gobble}
	
	\begin{titlepage}
		\centering
		\vspace*{4cm}
		
		{\Huge\bfseries 20 Lessons of Being 20\par}
		\vspace{1cm}
		
		{\Large dari manusia yang gagal\par}
		\vspace{3cm}
		
		{\Large Devangga\par}
		\vspace{0.5cm}
		{\large 2026\par}
		
		\vfill
	\end{titlepage}
	
	\cleardoublepage
	
	% =========================
	% FRONT MATTER (ROMAWI)
	% =========================
	\pagenumbering{roman}
	\setcounter{page}{1}
	
	% ----- HAK CIPTA -----
	\thispagestyle{empty}
	\vspace*{8cm}
	\begin{center}
		© 2026 Devangga\\
		Seluruh hak cipta dilindungi.
	\end{center}
	
	\cleardoublepage
	
	% ----- DAFTAR ISI -----
	\tableofcontents
	\cleardoublepage
	
	% ----- PENGANTAR -----
	\chapter*{Pengantar}
	\addcontentsline{toc}{chapter}{Pengantar}
	
	Dari antah berantah, dari daerah marginal, dari pengalaman pengalaman indah dan runyam aku menulis.
	
	Buku ini tak lain hanyalah pengingat dalam bentuk cerita yang ditujukan untuk penulis bahwa penulis. Sebelum ke 20 tahun penulis mendapat banyak pelajaran berharga diluar karir akademik.
	
	Teruntuk kamu yang belum mencapai umur 20 tahun, baca, ingat, dan bertidaklah sebelum difase ini, percayalah semakin cepat kamu bertindak, kehidupan 20 an mu akan jauh lebih mudah dari yang mungkin penulis sendiri rasakan, paling tidak buku ini membantumu sedikit menemukan tujuan, meski tidak dapat membanttumu menemukan arah kehidupan, tuhan menyertaimu. Mari kita jalani hidup sebaik mungkin.	\cleardoublepage
	
	% =========================
	% MAIN MATTER (ARABIC, MULAI 1)
	% =========================
	\pagenumbering{arabic}
	\setcounter{page}{1}
	\pagestyle{fancy}
	
	\chapter{Pelajaran Pertama: Kesadaran adalah kunci}
	Kita sebut dia ALex, pendek, kusam tak menarik perhatian. Kelas 8 saat itu ricuh sekali, ada seekor monyet bersekolah yang sedang menyatakan perasaannya pada dewi(Manusia paling cantik setelah ibu Alex). "Dewi mau kah kamu menjadi my pacar gua?" ucap Alex tak tau malu, "Sudah kubilang berapa kali si lex, gua ga tertarik sama elo, lagian gua nggak dibolehin pacaran sama bokap gua" ucap Dewi kepada monyet yang sedang mengungkapkan perasannya. Alex terdiam termenung dan tertegun malu karena baru saja ia di tolak oleh gadis yang ia sukai sendari kelas 7 dulu.
	
	"Ditolak berapa kalipun, kalau itu ditolak dewi, aku terima dengan lapang dada, meski sambil nangis di toilet." ucap Alex didalam toilet sambil mengusap air mata.
	
	\cleardoublepage
	
	\chapter{Pelajaran Kedua: Konsistensi Lebbih Penting dari Bakat}
	
	Tuliskan isi pelajaranmu di sini.
	
\end{document}
